% !TeX root = ./main.tex

\chapter{实验设计与结果分析}\label{ch:experiments}

本章系统评估多模态 GraphRAG 系统的问答性能，介绍实验设置、消融实验结果和错误分析。

\section{实验设置}\label{sec:exp-setup}

\subsection{数据集}

我们从 pdfQA/ClimRetrieve 文档集中选取 12 篇 ESG 可持续发展报告（涵盖 Microsoft、AT\&T、AstraZeneca、Boeing、Coca-Cola、Meta、Walmart 等多个行业），构建了包含 \textbf{120 道问答对}的自建测试集，覆盖六种问题类型：文本定位（30 题，25\%）、表格问答（27 题，22.5\%）、图表理解（17 题，14.2\%）、跨页延续（22 题，18.3\%）、多跳关系（12 题，10\%）和引用追踪（12 题，10\%）。每道问题标注了标准答案、问题类型、难度等级和证据页码。

\subsection{评测指标}

\begin{enumerate}[label=（\arabic*）]
\item \textbf{ANLS}：DocVQA 官方指标，归一化编辑距离（阈值 0.5）。
\item \textbf{Token F1}：基于字符 bigram 的 F1 分数，对中文友好。
\item \textbf{Heuristic Judge}：规则组合评分，EM=1.0、Containment=0.9、Token F1$\ge$0.8=0.8、Token F1$\ge$0.5=0.5、Token F1$\ge$0.3=0.25、其余 0。
\item \textbf{Evidence Page Recall}：预测页码与标准证据页码有交集的样本比例。
\end{enumerate}

所有指标均经 NFKC 规范化处理，多标准答案取下最大值。

\subsection{消融方案}

五组消融实验设计如下：

\begin{enumerate}[label=（\arabic*）]
\item \textbf{none}：该组不使用 FAISS 向量检索、异构文档图和页面图像证据，用于观察模型在缺少外部文档证据时的基线表现。
\item \textbf{vector\_only}：该组仅使用 FAISS 向量检索定位相关文本证据，用于检验纯语义召回在多页文档问答中的作用。
\item \textbf{graph\_only}：该组仅使用异构文档图进行证据扩展，用于检验结构关系和跨页关联对回答生成的贡献。
\item \textbf{no\_image}：该组同时使用 FAISS 向量检索和异构文档图，但不输入页面图像证据，用于分析视觉信息对表格、图表和版式相关问题的影响。
\item \textbf{full}：该组完整使用 FAISS 向量检索、异构文档图和页面图像证据，用于代表系统的完整多模态 GraphRAG 流程。
\end{enumerate}

所有实验使用相同的 120 道问题和 Qwen3-VL-8B-Instruct 模型，仅改变证据召回和输入配置，以保证不同组别之间具有可比性。

\section{实验结果}\label{sec:exp-results}

\subsection{整体性能}

五组消融实验的整体结果如图~\ref{fig:ablation-metrics} 和表~\ref{tab:ablation-results} 所示。完整系统在所有四项指标上均取得最优成绩。

\begin{figure}[htbp]
\centering
\includegraphics[width=\textwidth,keepaspectratio]{ablation_metrics.png}
\caption{五组消融实验四项核心指标对比}
\label{fig:ablation-metrics}
\end{figure}

\begin{table}[htbp]
\centering
\caption{五组消融实验整体指标数值}
\label{tab:ablation-results}
\begin{tabular}{@{}lrrrrr@{}}
\toprule
\textbf{指标} & \textbf{none} & \textbf{graph\_only} & \textbf{vector\_only} & \textbf{no\_image} & \textbf{full} \\
\midrule
ANLS & 0.0088 & 0.0044 & 0.0549 & 0.0790 & \textbf{0.0840} \\
Token F1 & 0.1924 & 0.4007 & 0.5704 & 0.6074 & \textbf{0.6185} \\
Heuristic Judge & 0.1042 & 0.3042 & 0.4408 & 0.4750 & \textbf{0.4813} \\
Evidence Page Recall & 0.0000 & 0.8917 & 0.8750 & 0.9167 & \textbf{0.9250} \\
\bottomrule
\end{tabular}
\end{table}

\subsection{消融分析}\label{sec:exp-ablation-analysis}

图~\ref{fig:waterfall} 展示了各组件的增量贡献（以 ANLS 为指标），核心发现如下：

\begin{figure}[htbp]
\centering
\includegraphics[width=0.85\textwidth,keepaspectratio]{waterfall.png}
\caption{组件增量贡献（ANLS 瀑布图）}
\label{fig:waterfall}
\end{figure}

\textbf{（1）FAISS 向量检索是最关键的组件。} vector\_only 的 ANLS（0.0549）是 graph\_only（0.0044）的 12.5 倍，Token F1 和 Heuristic Judge 也显著领先。仅有图谱关键词匹配无法提供足够的语义覆盖，FAISS 的稠密语义召回是不可替代的基础。

\textbf{（2）图谱扩展提供额外增益。} 对比 full 与 vector\_only，引入图谱后 ANLS 提升 53.1\%（0.0549 $\rightarrow$ 0.0840），Token F1 提升 8.4\%，Evidence Page Recall 提升 5.7\%。图谱邻域扩展增强了语义证据广度，提升了页面定位准确性。

\textbf{（3）页面图像有正向贡献但幅度较小。} 移除图像后 ANLS 下降 6.0\%，Token F1 下降 1.8\%，说明当前文档的文本抽取质量较高，图像主要对图表理解和视觉细节问题发挥作用。

\textbf{（4）graph\_only 的``高召回、低准确''现象值得关注。} graph\_only 的 Evidence Page Recall（0.8917）与 vector\_only（0.8750）相当，但 ANLS 仅为 0.0044。图谱可通过关键词找到正确页面，但从页面提取的文本证据质量不足，导致模型``看对了地方''却``答不对问题''。

\subsection{按问题类型分析}

图~\ref{fig:by-type} 和表~\ref{tab:by-type} 展示了完整系统在六种问题类型上的性能分布。text\_location 和 table\_qa 表现最好（ANLS 0.15 和 0.12，Heuristic Judge 约 0.50），FAISS 检索能准确定位单页内的文本和表格信息。cross\_page 和 chart\_understanding 居中，跨页证据合并和图表视觉推理仍有改进空间。multi\_hop 和 citation\_tracking 的 ANLS 为零，但 Token F1 约 0.59，说明回答语义上存在重叠——零 ANLS 主要因长答案多等价表述导致编辑距离评分失效。

\begin{figure}[htbp]
\centering
\includegraphics[width=\textwidth,keepaspectratio]{by_question_type.png}
\caption{完整系统按问题类型的 ANLS 与 Token F1 对比}
\label{fig:by-type}
\end{figure}

\begin{table}[htbp]
\centering
\caption{完整系统按问题类型的指标分布}
\label{tab:by-type}
\begin{tabular}{@{}lrrrr@{}}
\toprule
\textbf{问题类型} & \textbf{数量} & \textbf{ANLS} & \textbf{Token F1} & \textbf{Heuristic Judge} \\
\midrule
text\_location & 30 & \textbf{0.1535} & 0.6305 & 0.4767 \\
table\_qa & 27 & 0.1218 & \textbf{0.6461} & \textbf{0.5019} \\
chart\_understanding & 17 & 0.0989 & 0.6102 & 0.4765 \\
cross\_page & 22 & 0.0232 & 0.6018 & 0.4795 \\
multi\_hop & 12 & 0.0000 & 0.5898 & 0.4792 \\
citation\_tracking & 12 & 0.0000 & 0.5976 & 0.4583 \\
\bottomrule
\end{tabular}
\end{table}

\section{Token 成本分析}\label{sec:exp-token}

DocRAG 每次问答仅向 VLM 发送检索筛选后的相关页面和文本证据，而非整个文档。为了量化这一设计对调用成本的影响，本文将 DocRAG 问答流程与直接上传目标 PDF 的全量问答方式进行对比。两种方式使用相同的问题集合，统计模型调用中的输入 token、输出 token 和总 token。实验共选取 20 道真实 PDF 问题，其中 13 道问题在两种方式下均成功返回结果，因此表~\ref{tab:token-cost} 以这 13 个有效配对样本作为统计口径。

\begin{table}[htbp]
\centering
\caption{DocRAG 与全量上传方式的 Token 成本对比}
\label{tab:token-cost}
\begin{tabular}{lccc}
\toprule
统计项 & DocRAG & 全量上传 PDF & 相对全量上传变化 \\
\midrule
有效配对样本数 & 13 & 13 & -- \\
输入 Token 总数 & 29,071 & 547,890 & 减少 94.69\% \\
总 Token 数 & 29,097 & 547,903 & 减少 94.69\% \\
平均输入 Token & 2,236 & 42,145 & 减少 94.69\% \\
平均总 Token & 2,238 & 42,146 & 减少 94.69\% \\
\bottomrule
\end{tabular}
\end{table}

从表~\ref{tab:token-cost} 可以看出，全量上传方式在每次问答中都需要重新处理完整 PDF，因此平均输入 token 达到 42,145；DocRAG 则通过检索和图谱扩展先定位少量相关证据，平均输入 token 降至 2,236，仅为全量上传方式的约 5.31\%。换言之，DocRAG 在输入侧节省了 94.69\% 的 token。输出 token 由于本实验要求模型给出简短答案，两种方式的数值都很小，且 DocRAG 略高于全量上传方式，因此不构成主要成本差异。总体来看，token 成本优势主要来自输入侧压缩，说明先检索再生成的流程能够在保持证据可追溯的同时显著降低长 PDF 问答的调用开销。

图~\ref{fig:token} 进一步直观展示了输入 token 节省比例。这一优势在多轮对话和批量评测场景下尤为显著，因为全量上传方式会在每轮问答中重复传入整篇文档，而 DocRAG 只需传入与当前问题相关的局部证据。

\begin{figure}[htbp]
\centering
\includegraphics[width=0.55\textwidth,keepaspectratio]{token_savings.png}
\caption{DocRAG vs 全量上传：估算输入 Token 节省率}
\label{fig:token}
\end{figure}

\section{错误分析}\label{sec:exp-error}

通过人工审查低分样本，归纳四类主要错误模式：

\begin{enumerate}[label=（\arabic*）]
\item \textbf{检索跨度不足}：FAISS 召回可能仅命中跨页起止页中的一页，continuation 边未将另一端文本引入证据集。
\item \textbf{答案格式冗长}：VLM 即使 short 模式也输出大量铺垫，导致与简洁标准答案的 ANLS/EM 偏低。
\item \textbf{图表视觉推理脱节}：模型有时仅依赖文本形式的图表描述，忽略页面原图中的视觉细节（趋势线斜率、柱状图相对高度）。
\item \textbf{多跳证据链断裂}：multi\_hop 问题需在不同节点间建立推理链，当中间节点的语义关系边存在噪声或缺失时，证据链中断。
\end{enumerate}

\section{人工评分交叉验证}\label{sec:exp-human}

自动指标（尤其是 ANLS 和 Exact Match）对长答案的语义等价不敏感：当模型给出内容正确但表述不同的回答时，编辑距离类指标往往给出低分。为交叉验证自动指标的可靠性，本节从测试集中选取了 4 篇文档（Meta 2023 Sustainability Report、Veolia ESG Report 2023、Walmart ESG Highlights 2023、Westpac 2023 Climate Report），对其共计 40 道问题进行人工三档评分，并与 Heuristic Judge 进行对比分析。

\subsection{评分标准}

人工评分采用三档制，以标准答案的任意变体为基准，不要求逐字匹配：

\begin{enumerate}[label=（\arabic*）]
\item \textbf{正确（1.0）}：核心事实与标准答案一致，关键数据/结论正确，无重大遗漏。
\item \textbf{部分正确（0.5）}：触及核心但存在关键数据偏差、遗漏重要子问题、部分信息不准确或回答过于泛化。
\item \textbf{错误（0.0）}：关键事实错误、答非所问、或模型表示``证据不足''但标准答案有明确信息。
\end{enumerate}

多子问题等权重独立判断；模型给出额外合理信息不扣分。加权准确率定义为： $$\text{加权准确率} = (N_{\text{正确}} \times 1.0 + N_{\text{部分正确}} \times 0.5) / 40$$

\subsection{整体结果}

40 道问题的人工评分整体分布如表~\ref{tab:human-eval-overall} 所示。

\begin{table}[htbp]
\centering
\caption{人工评分整体分布（40 题）}
\label{tab:human-eval-overall}
\begin{tabular}{@{}lrr@{}}
\toprule
\textbf{评分等级} & \textbf{数量} & \textbf{占比} \\
\midrule
正确（1.0）& 19 & 47.5\% \\
部分正确（0.5）& 15 & 37.5\% \\
错误（0.0）& 6 & 15.0\% \\
\midrule
加权准确率 & & \textbf{0.6625} \\
\bottomrule
\end{tabular}
\end{table}

加权准确率为 66.3\%，显著高于 Heuristic Judge 对同 40 题的平均分（0.496）。这说明模型回答的实质语义质量优于编辑距离类指标所反映的水平，大量``部分正确''回答被自动指标系统性低估。

\subsection{按文档分析}

四篇文档的人工加权准确率如表~\ref{tab:human-eval-doc} 所示。

\begin{table}[htbp]
\centering
\caption{人工评分与 Heuristic Judge 按文档对比}
\label{tab:human-eval-doc}
\begin{tabular}{@{}lrrr@{}}
\toprule
\textbf{文档} & \textbf{人工加权准确率} & \textbf{Heuristic Judge} & \textbf{差距} \\
\midrule
Veolia ESG Report 2023 & 0.750 & 0.500 & +0.250 \\
Meta 2023 Sustainability Report & 0.700 & 0.530 & +0.170 \\
Walmart ESG Highlights 2023 & 0.600 & 0.425 & +0.175 \\
Westpac 2023 Climate Report & 0.600 & 0.530 & +0.070 \\
\midrule
平均 & 0.663 & 0.496 & +0.167 \\
\bottomrule
\end{tabular}
\end{table}

Veolia 表现最好（0.750），Walmart 和 Westpac 相对较弱（0.600）。值得注意的是，Westpac 的 Heuristic Judge 评分（0.530）与 Veolia（0.500）相当甚至略高，但人工评分差距达 0.150，再次说明自动指标可能高估部分文档的表现。

\subsection{按题型分析}

各题型的表现差异明显（表~\ref{tab:human-eval-type}）。

\begin{table}[htbp]
\centering
\caption{人工评分按题型分布}
\label{tab:human-eval-type}
\begin{tabular}{@{}lrrrr@{}}
\toprule
\textbf{题型} & \textbf{数量} & \textbf{加权准确率} & \textbf{Heuristic Judge} & \textbf{差距} \\
\midrule
text\_location & 12 & 0.792 & 0.500 & +0.292 \\
table\_qa & 9 & 0.722 & 0.506 & +0.217 \\
chart\_understanding & 4 & 0.625 & 0.500 & +0.125 \\
cross\_page & 7 & 0.643 & 0.471 & +0.172 \\
citation\_tracking & 4 & 0.375 & 0.500 & $-$0.125 \\
multi\_hop & 4 & 0.125 & 0.500 & $-$0.375 \\
\bottomrule
\end{tabular}
\end{table}

text\_location 和 table\_qa 表现最好（加权准确率 > 0.7），这两类问题对单页文本和表格信息的定位需求明确，FAISS 检索能有效覆盖。multi\_hop 表现最弱（4 题中 3 题得 0 分，加权准确率仅 0.125），主要原因是中间节点的语义关系边存在噪声或缺失，导致多跳推理链条中断。citation\_tracking 次弱（0.375），模型在给定页码时频频出错。

一个令人警惕的现象是：multi\_hop 和 citation\_tracking 的人工加权准确率远低于 Heuristic Judge（差距分别为 $-$0.375 和 $-$0.125），这意味着自动指标在这两类题型上存在严重的假阳性问题——模型输出的文本因包含实体名称和数值关键词而获得较高自动分，但实际回答的关键事实是错误的。

\subsection{自动指标可靠性分析}

将人工评分与 Heuristic Judge 逐题配对进行 Spearman 秩相关性分析，系数仅为 \textbf{0.387}（$p < 0.05$），属于弱正相关。进一步识别出以下典型不一致案例：

\begin{enumerate}[label=（\arabic*）]
\item \textbf{假阳性（5 例）}：Heuristic Judge $\ge 0.5$ 但人工评分为 0.0。这 5 例全部集中在 multi\_hop 和 citation\_tracking 题型，共同特征是模型输出包含与标准答案相同的实体名称和数值，但关键关系或归因完全错误，自动指标仅凭表面词级匹配无法识别。
\item \textbf{假阴性（1 例）}：Heuristic Judge = 0.25 但人工评分为 1.0（Walmart 员工构成数据题）。模型回答数据完全正确，但因表述顺序与标准答案不同导致 Token F1 下降，进而拉低 Heuristic Judge。
\end{enumerate}

综合以上分析，本节的人工评分交叉验证得出三点主要结论：第一，完整系统在人工三档评分下的加权准确率达到 66.3\%，高于自动指标所反映的水平，系统具有实际可用性；第二，自动指标（尤其是 Heuristic Judge 和 ANLS）对 multi\_hop 和 citation\_tracking 类长答案存在系统性高估或低估，不应单独作为评测依据；第三，在资源允许的情况下，建议将人工抽检作为自动评测的必要补充。

\section{实验小结}\label{sec:exp-summary}

本章通过五组消融实验验证了``向量检索 + 图谱扩展 + 多模态融合''路线的有效性：完整系统在全部指标上最优，FAISS 语义召回是最关键的单一组件，图谱和图像提供额外增益，页面级证据定位达 92.5\% 召回率。Token 分析表明 DocRAG 相比全量上传可节省约 94\% 的输入 token。
